% !PRINTED 332
\weberpage{332}
From the expression (1) we can derive a \gesperrt{canonical representation}
 for the cubic form $f$:
We denote by $\varepsilon$ an imaginary third root of unity
and introduce two variables $z_1, z_2$, which are defined by the equations

\weq{2}{\begin{aligned} a_1\,x_1 &= \varepsilon\,z_1 + \varepsilon^2\,z_2 - \tfrac{1}{3}\,a_3\,x_3 \\ a_2\,x_2 &= \varepsilon^2\,z_1 + \varepsilon\,z_2 - \tfrac{1}{3}\,a_3\,x_3 \end{aligned}}

from which it further follows:

\weq{3}{-\,(a_1\,x_1 + a_2\,x_2 + a_3\,x_3) = z_1 + z_2 - \tfrac{1}{3}\,a_3\,x_3.}

The variables $z_1, z_2$ are determined by (2) as linear functions
of $x_1, x_2, x_3$, and the linear functions $z_1, z_2, x_3$
are linearly independent as functions of $x_1, x_2, x_3$. If we
introduce (2) and (3) into (1), we obtain

\[
-\,a_1\,a_2\,f = z_1^3 + z_2^3 - \left(a_1\,a_2\,h + \frac{a_3^3}{27}\right)x_3^3 + a_3\,z_1\,z_2\,x_3.
\]

One gives this expression a more symmetric shape
 by introducing for $z_1, z_2, x_3$ three new linear functions
$y_1, y_2, y_3$, which differ from these only by constant factors
:

\[
f = \varphi\,(y_1,y_2,y_3) = \alpha_1\,y_1^3 + \alpha_2\,y_2^3 + \alpha_3\,y_3^3 + 6\,m\,y_1\,y_2\,y_3,
\]

where $\alpha_1, \alpha_2, \alpha_3, m$ denote constants.

If no singular point is present, the
coefficients $\alpha_1, \alpha_2, \alpha_3$ cannot vanish, and one can,
without restriction of generality, take them equal to one another, say $=h$.
One could even set them $=1$, which however, in order to
maintain homogeneity, is for the present better not done.
We thus have the \gesperrt{canonical form} for the ternary
cubic form:

\weq{4}{\varphi\,(y_1,y_2,y_3) = h\,(y_1^3+y_2^3+y_3^3) + 6\,m\,y_1\,y_2\,y_3.}

According to our point of departure, the line $x_3$, or what
is the same thing, $y_3$, was the line joining three inflexion points.
But there are twelve such lines, and we can therefore choose the
line $y_3$ in twelve ways. If, however, we put $y_1$ or $y_2$ in
place of $y_3$, we obtain the same representation in the
canonical form, and there are thus only four essentially different
 kinds of this representation (apart from the arbitrary
 $h$), i.\,e. four ways of determining the lines $y_1, y_2, y_3$. We
therefore have the theorem:
%CONT>>
